% Fichier complété par tools/complete_missing_corrections.py.
% Source : DYN/DYN-06-PFD/09_RT_RSG/09_RT_RSG.tex
% Statut : rédigé (2 question(s)).
\begin{corrigebox}[Corrigé rédigé]
\CorrigeQuestion{1}
Sur $\vec i_1$, la résultante dynamique de 2 donne
                $\sum F_{i_1}^{\rm ext}=m_2\vec\Gamma(G_2,2/0)\cdot\vec i_1$.
                En utilisant l'accélération cinématique et la contrainte de roulement, on
                obtient l'équation de translation de $\lambda$.
\CorrigeQuestion{2}
Au point de contact $I$, le moment dynamique de 1+2 projeté sur
                $\vec k_0$ élimine les réactions du sol :
                $\sum M_{I,k_0}^{\rm ext}=\delta_{I,k_0}(1+2/0)$. Cette équation donne la
                loi en $\theta$ couplée à $\lambda$.
\end{corrigebox}
